% operationqueue-deep.tex — the Operation state machine, KVO-compliant async Operation,
% dependencies, maxConcurrentOperationCount, QoS vs queuePriority, cooperative
% cancellation, BlockOperation, OperationQueue.main, Operation vs GCD vs async/await,
% common bugs.
% Sources (checked 2026-10-05):
%   developer.apple.com/documentation/foundation/operation (state KVO keys; a queue ignores
%     isAsynchronous and always calls start() on a separate thread; cancel() — unfinished
%     dependencies ignored; completionBlock set to nil once it starts, iOS 8+; circular
%     dependencies deadlock; QueuePriority raw values)
%   developer.apple.com/documentation/foundation/operationqueue (defaultMaxConcurrentOperationCount
%     = -1; addBarrierBlock and progress, iOS 13; underlyingQueue; isSuspended)
%   developer.apple.com/documentation/foundation/qualityofservice (raw values)
%   Concurrency Programming Guide, "Operation Queues" (async operation KVO)
% @labels: area=ios-platform kind=api platform=apple topic=concurrency discussion=39
% @tags: operation, operationqueue, async-operation, kvo, isfinished, adddependency, maxconcurrentoperationcount, cancellation, queuepriority, qos, blockoperation, gcd
\documentclass[8pt]{extarticle}
\usepackage{printup-sheet}
\usepackage{array}

\lstdefinelanguage{SwiftSheet}{
  morekeywords={protocol,class,final,struct,enum,func,var,let,weak,init,override,
    if,else,return,guard,self,nil,try,await,async,throws,private,some,defer,
    true,false,AnyObject,Void,String,Bool},
  sensitive=true, morecomment=[l]{//}, morecomment=[s]{/*}{*/}, morestring=[b]"}

\tikzset{
  st/.style={box, font=\scriptsize, inner sep=1.5pt, minimum height=8mm, minimum width=15mm},
  lbl/.style={font=\tiny, text=black!80, inner sep=1pt, align=center},
  op/.style={box, font=\scriptsize, inner sep=1.5pt, minimum height=5.5mm, minimum width=14mm},
  cx/.style={op, draw=sheetRed, fill=sheetRed!8},
}
\newcolumntype{L}[1]{>{\raggedright\arraybackslash}p{#1}}

\begin{document}

\sheettitle{Operation \& OperationQueue}{ios-platform · folio}

\oneliner{An \texttt{Operation} is a \textbf{single-shot}, KVO-observable unit of work;
an \texttt{OperationQueue} starts it when \texttt{isReady} (every dependency
\texttt{isFinished}), caps parallelism with \texttt{maxConcurrentOperationCount}, and drops it
when it posts \texttt{isFinished}. \textbf{Cancellation is a flag}, dependencies are
\textbf{ordering only} (no data), and an async subclass must drive its own state
\textbf{with KVO} — or it never finishes.}

\vspace{2pt}
\noindent\begin{tikzpicture}[sheet]
  % ── state machine ──
  \node[font=\bfseries\small, anchor=west] at (-0.1,2.85) {State machine (KVO keys the queue watches)};
  \node[st, draw=sheetGrey, fill=black!5] (p) at (0.65,2.0) {pending\\deps unfinished};
  \node[st] (r) at (3.1,2.0) {\texttt{isReady}};
  \node[st, draw=sheetOrange, fill=sheetOrange!10] (e) at (5.8,2.0) {\texttt{isExecuting}};
  \node[st, draw=sheetGreen!70!black, fill=sheetGreen!10] (f) at (8.95,2.0) {\texttt{isFinished}\\left the queue};
  \draw[flow] (p) -- node[lbl, above]{deps all}  node[lbl, below]{finished} (r);
  \draw[flow] (r) -- node[lbl, above]{queue calls} node[lbl, below]{\texttt{start()}} (e);
  \draw[flow] (e) -- node[lbl, above]{\texttt{main()} returns} node[lbl, below]{or \texttt{finish()}} (f);
  \draw[->, thick, sheetRed, rounded corners=4pt] (p.south) -- ++(0,-0.5) -| (f.south);
  \node[lbl, text=sheetRed, fill=white] at (4.9,1.1)
    {\texttt{cancel()} before start: deps \textbf{ignored} → ready → \texttt{start()} sees the flag → finished; \texttt{main()} never runs};
  \node[lbl, anchor=west, align=left, text=sheetRed] at (-0.1,0.55)
    {\texttt{isCancelled} is an orthogonal \textbf{flag}, settable in any state. While executing
     it changes nothing\\unless \emph{your} code checks it.};
  \node[lbl, anchor=west, align=left, text=sheetBrown] at (-0.1,0.0)
    {async subclass without \texttt{willChange/didChange} for \texttt{"isFinished"}: the queue
     never sees the end →\\dependents never start and the concurrency slot is never freed.};
  % ── dependency graph ──
  \draw[sheetGrey!40] (10.05,2.95) -- (10.05,-0.25);
  \node[font=\bfseries\small, anchor=west] at (10.15,2.85) {Dependencies — a graph across queues};
  \node[lbl, anchor=west, text=sheetBlue] at (10.15,2.45) {\texttt{net} queue, \texttt{maxConcurrentOperationCount = 2}};
  \node[op] (d1) at (11.1,1.9) {download A};
  \node[cx] (d2) at (11.1,1.05) {download B\\\textbf{cancelled}};
  \node[op] (dc) at (13.1,1.5) {decode};
  \node[op, draw=sheetGreen!70!black, fill=sheetGreen!10] (ui) at (15.35,1.5) {show\\on \texttt{.main}};
  \draw[flow] (d1) -- (dc); \draw[flow, sheetRed] (d2) -- (dc); \draw[flow] (dc) -- (ui);
  \node[lbl, text=sheetRed, align=left, anchor=west] at (10.15,0.35)
    {a cancelled op still \textbf{finishes} → \texttt{decode} runs anyway:\\
     check \texttt{dependencies} for \texttt{isCancelled} / an error};
  \node[lbl, anchor=west] at (10.15,-0.1) {\texttt{decode.addDependency(dA)} — arrow = ``must finish before''};
\end{tikzpicture}

\begin{multicols}{2}
\raggedright

\section{How it works}
\begin{itemize}
  \item \textbf{Sync op}: override \texttt{main()}; finished when it returns (async work
        started inside is \emph{not} awaited). \textbf{Async op}: override
        \texttt{start()} (never \texttt{super.start()}), \texttt{isAsynchronous = true},
        own \texttt{isExecuting}/\texttt{isFinished} behind a lock, KVO for both, finish
        exactly once — also when cancelled before starting.
  \item A queue \textbf{always} calls \texttt{start()} on a separate thread and ignores
        \texttt{isAsynchronous}; it matters only when you call \texttt{start()} yourself.
  \item \texttt{b.addDependency(a)}: works across queues; ``finished'' includes cancelled; a
        cycle never becomes ready (silent hang). Results: \texttt{b} reads
        \texttt{a.output} (or an adapter op copies it).
  \item \texttt{maxConcurrentOperationCount} defaults to −1 (system decides);
        \texttt{1} = serial, but order = readiness → \texttt{queuePriority} → insertion.
  \item \texttt{qualityOfService} (\texttt{.userInteractive … .background}) sets the
        thread's priority and energy class; \texttt{queuePriority}
        (\texttt{.veryLow … .veryHigh}) only reorders \emph{ready} ops in \emph{one} queue.
  \item \texttt{BlockOperation}: its blocks run concurrently, finished when all return.
        \texttt{addBarrierBlock} (iOS 13) waits for everything queued before it.
        \texttt{isSuspended} stops \emph{starting} ops. \texttt{OperationQueue.main} =
        serial, main thread; \texttt{underlyingQueue} targets a \texttt{DispatchQueue}.
  \item \texttt{completionBlock} runs after finish on an \textbf{arbitrary thread}, and is
        set to \texttt{nil} once it starts (iOS 8+).
\end{itemize}

\section{Example — KVO-correct async Operation}
\begin{lstlisting}[language=SwiftSheet]
class AsyncOperation: Operation {
  private let lock = NSLock()
  private var _exec = false, _done = false
  private func locked<T>(_ f: () -> T) -> T {
    lock.lock(); defer { lock.unlock() }; return f() }
  override var isAsynchronous: Bool { true }
  override var isExecuting: Bool { locked { _exec } }
  override var isFinished: Bool { locked { _done } }
  override func start() {              // never super.start()
    if isCancelled { finish(); return } // still must finish
    willChangeValue(forKey: "isExecuting")
    locked { _exec = true }
    didChangeValue(forKey: "isExecuting")
    run { self.finish() } }            // subclass: async work
  func run(_ done: @escaping () -> Void) { done() }
  final func finish() {                // exactly once
    let keys = ["isExecuting", "isFinished"]
    keys.forEach(willChangeValue(forKey:))
    locked { _exec = false; _done = true }
    keys.forEach(didChangeValue(forKey:)) }
}
\end{lstlisting}

\columnbreak

\section{Operation vs GCD vs async/await}
{\footnotesize\setlength{\tabcolsep}{2.5pt}
\begin{tabular}{@{}L{14mm}L{21mm}L{19mm}L{20mm}@{}}
\toprule
\textbf{Need} & \textbf{Operation} & \textbf{GCD} & \textbf{Swift concurrency} \\
\midrule
dependency graph & \texttt{addDependency}, cross-queue & group \texttt{notify}, by hand & \texttt{async let}, ordered \texttt{await} \\
cancel a graph & \texttt{cancelAll…()} sets flags & work item \texttt{cancel()}: only if not started & \texttt{task.cancel()} reaches child tasks \\
cap parallelism & \texttt{maxConcurrent…} & semaphore / serial queue & task group: seed $n$, add one per finish \\
observe state & KVO, \texttt{progress} (iOS 13) & — & — \\
delay & — & \texttt{asyncAfter} & \texttt{Task.sleep(for:)} \\
\bottomrule
\end{tabular}}
{\footnotesize Operation still wins for: a \textbf{cancellable graph} across queues, a hard
concurrency cap, KVO/\texttt{Progress} UI, subclassable reusable \emph{types}, legacy code.
New code on async/await: task groups + \texttt{Task.checkCancellation()}.}

\section{Traps}
\begin{itemize}
  \trap{Async op without the manual KVO → stuck ``executing'' forever.}
  \trap{\texttt{cancel()} does not stop a running op or its \texttt{URLSessionTask} —
        check \texttt{isCancelled} and cancel the task yourself.}
  \trap{Cancelled dependency = finished → the dependent still runs.}
  \trap{\texttt{op.completionBlock = \{ use(op.result) \}} captures \texttt{op} strongly —
        a cycle until it runs; use \texttt{[unowned op]}. Hop to main for UI.}
  \trap{Adding the same instance twice / after it finished →
        \texttt{NSInvalidArgumentException}. Single-shot.}
  \trap{\texttt{waitUntilAllOperationsAreFinished()} / \texttt{waitUntilFinished: true} on
        main → hang, watchdog \texttt{0x8badf00d}.}
\end{itemize}

\section{The two priorities, side by side}
{\footnotesize
\begin{tabular}{@{}l r l r@{}}
\toprule
\texttt{qualityOfService} & raw & \texttt{queuePriority} & raw \\ \midrule
\texttt{.userInteractive} & 0x21 & \texttt{.veryHigh} & 8 \\
\texttt{.userInitiated} & 0x19 & \texttt{.high} & 4 \\
\texttt{.default} & $-1$ & \texttt{.normal} & 0 \\
\texttt{.utility} & 0x11 & \texttt{.low} & $-4$ \\
\texttt{.background} & 0x09 & \texttt{.veryLow} & $-8$ \\
\bottomrule
\end{tabular}}\\[1pt]
{\scriptsize QoS = which CPU/energy class the \emph{thread} gets (across the whole system);
queue priority = which \emph{ready} op this queue starts first. Neither beats a dependency.}

\section{Remember}
\textbf{Ready → Executing → Finished; cancel is only a flag; no KVO, no finish.}


\end{multicols}

\noindent{\footnotesize\color{sheetGrey}\textit{Related:} concurrency (GCD, actors) ·
concurrency-patterns (bounded parallelism) · async-sequences-streams · KVO
(observation-and-combine) · background-execution · crashes-symbolication (0x8badf00d)}

\end{document}
